\documentclass[11pt]{article}
\pdfoutput=1
\usepackage{amsmath}
\usepackage{hyperref}
\usepackage{geometry}
\usepackage{listings}
\usepackage{color}
\usepackage{booktabs}
\usepackage{longtable}
\usepackage{array}
\usepackage{parskip}
\usepackage{fancyhdr}
\usepackage{titlesec}
\usepackage{mdframed}

\geometry{
  letterpaper,
  top=25mm, bottom=25mm,
  left=25mm, right=25mm
}

% ---------------------------------------------------------------
% Code listing style
% ---------------------------------------------------------------
\definecolor{codebg}{rgb}{0.95,0.95,0.95}
\definecolor{codeframe}{rgb}{0.7,0.7,0.7}
\definecolor{commentcolor}{rgb}{0.4,0.4,0.4}
\definecolor{keywordcolor}{rgb}{0.1,0.3,0.7}

\lstset{
  backgroundcolor=\color{codebg},
  basicstyle=\small\ttfamily,
  breaklines=true,
  frame=single,
  rulecolor=\color{codeframe},
  tabsize=4,
  columns=flexible,
  keepspaces=true,
  showstringspaces=false,
  commentstyle=\color{commentcolor}\itshape,
  keywordstyle=\color{keywordcolor},
  numbers=none,
  xleftmargin=6pt,
  xrightmargin=6pt,
  aboveskip=8pt,
  belowskip=8pt
}

% Convenience macro for inline code
\newcommand{\code}[1]{\texttt{#1}}

% ---------------------------------------------------------------
% Note / Warning boxes
% ---------------------------------------------------------------
\newmdenv[
%  backgroundcolor=yellow!12,
%  linecolor=orange!60,
  linewidth=1pt,
  leftmargin=0pt, rightmargin=0pt,
  innerleftmargin=8pt, innerrightmargin=8pt,
  innertopmargin=6pt, innerbottommargin=6pt,
  skipabove=8pt, skipbelow=8pt
]{notebox}

\newmdenv[
  backgroundcolor=red!8,
%  linecolor=red!60,
  linewidth=1pt,
  leftmargin=0pt, rightmargin=0pt,
  innerleftmargin=8pt, innerrightmargin=8pt,
  innertopmargin=6pt, innerbottommargin=6pt,
  skipabove=8pt, skipbelow=8pt
]{warnbox}

% ---------------------------------------------------------------
% Headers / footers
% ---------------------------------------------------------------
\pagestyle{fancy}
\fancyhf{}
\fancyhead[L]{\small Magic2 User Guide}
\fancyhead[R]{\small v1.0}
\fancyfoot[C]{\small\thepage}
\renewcommand{\headrulewidth}{0.4pt}

% ---------------------------------------------------------------
% Section formatting
% ---------------------------------------------------------------
\titleformat{\section}{\large\bfseries}{{\thesection}}{0.8em}{}[\vspace{2pt}\hrule\vspace{4pt}]
\titleformat{\subsection}{\normalsize\bfseries}{{\thesubsection}}{0.8em}{}
\titleformat{\subsubsection}{\normalsize\itshape\bfseries}{{\thesubsubsection}}{0.8em}{}

% ---------------------------------------------------------------
\begin{document}

% ---------------------------------------------------------------
% Title page
% ---------------------------------------------------------------
\begin{titlepage}
\vspace*{20mm}
{\Huge \textbf{Magic2}}\\[4mm]
{\Large Deep RNA-Sequencing Analyser}\\[2mm]
{\large User Guide}
\vspace{8mm}
\hrule
\vspace{8mm}
{\large Jean ~Thierry-Mieg, Danielle ~Thierry-Mieg, Greg ~Boratyn}\\[2mm]
National Library of Medicine, National Institutes of Health\\
8600 Rockville Pike, Bethesda MD 20894, USA\\[2mm]
\href{mailto:mieg@ncbi.nlm.nih.gov}{\texttt{mieg@ncbi.nlm.nih.gov}}
\href{mailto:boratyng@ncbi.nlm.nih.gov}{\texttt{boratyng@ncbi.nlm.nih.gov}}
\vspace{8mm}

\noindent Source code: \href{https://github.com/NLM-DIR/Magic2}{\texttt{https://github.com/NLM-DIR/Magic2}}

\vfill
\small National Library of Medicine $\cdot$ National Institutes of Health
\end{titlepage}

% ---------------------------------------------------------------
\newpage
\tableofcontents
\newpage

% ===============================================================
\section{Overview}
% ===============================================================

Magic2 is an accurate, sensitive RNA deep-sequencing aligner designed for
both short reads (Illumina paired-end) and long reads (PacBio, Nanopore,
Roche).  In a single pass it produces alignment files, gene expression
counts, intron support tables, poly-A addition sites, coverage plots, and
quality-control metrics.

\textbf{Key capabilities:}
\begin{itemize}
  \item Aligns reads against a raw genome or an annotated genome (GTF/GFF).
  \item Reads local FASTA/FASTQ files (gzip-compressed or plain) or streams
        data on the fly from the NCBI Short Read Archive (SRA).
  \item Automatically recognises paired-end vs.\ single-end reads, adaptor
        sequences, sequencing technology, and read length---no strategy
        parameters required.
  \item Adapts its parallelism to the available hardware; offloads work to
        a GPU when one is present.
  \item Exports alignments in BAM, Magic \code{.hits}, or tabular format.
\end{itemize}

\begin{notebox}
\textbf{Note:} Magic2 runs a dry run on the first 10~Megabases to detect
read characteristics before committing to a full alignment strategy.
You do not need to specify whether reads are DNA or RNA, short or long,
or whether they contain adaptors.
\end{notebox}

% ===============================================================
\section{System Requirements}
% ===============================================================

\begin{center}
\begin{tabular}{ll}
\toprule
\textbf{Component} & \textbf{Minimum / Recommended} \\
\midrule
Operating system & Linux (x86-64) or macOS \\
RAM              & 32 GB minimum; 64+ GB recommended for human genome \\
CPU cores        & 4 minimum; performance scales linearly with core count \\
Disk             & Varies with genome size; \(\sim\)30 GB for a human index \\
GPU (optional)   & Any CUDA-capable GPU accelerates seed extension \\
Network          & Fast connection to NCBI recommended for SRA streaming \\
\bottomrule
\end{tabular}
\end{center}

% ===============================================================
\section{Installation}
% ===============================================================

\subsection{Using pre-built executables}

Pre-compiled binaries for Linux and macOS are available from:
\begin{lstlisting}
https://github.com/NLM-DIR/Magic2/releases
\end{lstlisting}

Download the archive for your platform, extract it, and add the
\code{bin/} directory to your \code{PATH}, or copy the \code{magic2}
executable to \code{/usr/local/bin}.

To verify the installation:
\begin{lstlisting}
magic2 --help
\end{lstlisting}

\subsection{Building from source}

The source code is available on GitHub under an open-source licence.

\begin{lstlisting}
git clone git@github.com:NLM-DIR/Magic2
\end{lstlisting}

Then compile:
\begin{lstlisting}
tcsh
cd Magic2
setenv ACEDB_MACHINE LINUX_4_OPT
make libs
cd wsa
make
\end{lstlisting}

The executable is created at \code{Magic2/bin.LINUX\_4\_OPT/magic2}.
Add this directory to your \code{PATH}:
\begin{lstlisting}
export PATH=$PATH:/path/to/Magic2/bin.LINUX_4_OPT
\end{lstlisting}

\begin{notebox}
\textbf{Note:} Locally compiled binaries may outperform the pre-built
releases because the compiler can target the specific instruction sets
of your hardware.
\end{notebox}

% ===============================================================
\section{Quick Start}
% ===============================================================

This section shows three minimal end-to-end examples to get you started
immediately.  Full details for each step are in the sections that follow.

\subsection{Example 1: Align reads against a raw genome}

\begin{lstlisting}
# Step 1: Build the index (one time only)
magic2 --createIndex IDX.genome -t genome.fasta.gz

# Step 2: Align reads
magic2 -x IDX.genome -i reads.fasta.gz -o results/
\end{lstlisting}

\code{genome.fasta.gz} must contain all chromosomes in a single FASTA
file.  Results appear in the \code{results/} directory.

\subsection{Example 2: Align reads against an annotated genome}

\begin{lstlisting}
# Step 1: Create a target configuration file  (tConfig)
#   Format: one tab-separated type/filename pair per line
#   G = genomic sequence   I = annotation (GTF/GFF)
#   M = mitochondria       R = ribosomal RNA
cat > human.tConfig << 'EOF'
G	T2T.chromosomes.fasta.gz
M	T2T.mitochondria.fasta.gz
R	T2T.rRNA.fasta.gz
I	T2T.annotation.gtf.gz
EOF

# Step 2: Build the annotated index
magic2 --createIndex IDX.T2T -T human.tConfig

# Step 3: Align and produce the full output suite
magic2 -x IDX.T2T -i reads.fastq.gz -o results/ --full
\end{lstlisting}

\subsection{Example 3: Stream and align SRA data}

\begin{lstlisting}
# Align three SRA runs on the fly, no local download required
magic2 -x IDX.T2T --srr SRR1234,SRR1235,SRR1236 -o results/ --full
\end{lstlisting}

% ===============================================================
\section{Input Data}
% ===============================================================

\subsection{Local FASTA/FASTQ files}

Magic2 accepts both FASTA and FASTQ files, gzip-compressed or
uncompressed.  The format is detected automatically.

\textbf{Single-end reads:}
\begin{lstlisting}
magic2 -x IDX -i reads.fasta.gz -o out/
\end{lstlisting}

\textbf{Paired-end reads in two separate files} (joined with \code{+}):
\begin{lstlisting}
magic2 -x IDX -i reads_R1.fastq.gz+reads_R2.fastq.gz -o out/
\end{lstlisting}

\textbf{Paired-end reads in interleaved format:}\\
The file must be named with the \code{.sample\_12.} suffix:
\begin{lstlisting}
magic2 -x IDX -i sample.sample_12.fastq.gz -o out/
\end{lstlisting}

In interleaved format, each pair occupies four successive lines and read
identifiers end in \code{/1} and \code{/2}.

\subsection{Downloading SRA runs}

Use \code{--sraDownload} to download one or more SRA runs as local files
before alignment.  The files are placed in the \code{./SRA/} subdirectory.

\textbf{Basic download (interleaved FASTA):}
\begin{lstlisting}
magic2 --sraDownload SRR24511885
# Output: SRA/SRR24511885.sample_12.fasta.gz
\end{lstlisting}

\textbf{Download as split FASTQ files, limiting to 2 GB:}
\begin{lstlisting}
magic2 --sraDownload SRR24511885 --fastq --split_pairs --maxGB 2
# Output: SRA/SRR24511885_R1.fastq.gz
#         SRA/SRR24511885_R2.fastq.gz
\end{lstlisting}

\textbf{Full syntax:}
\begin{lstlisting}
magic2 --sraDownload SRR1,SRR2,... [--fastq] [--split_pairs] [--maxGB <n>]
\end{lstlisting}

\begin{center}
\begin{tabular}{lp{9cm}}
\toprule
\textbf{Option} & \textbf{Description} \\
\midrule
\code{--fastq}       & Download in FASTQ format (preserves quality scores).
                        Default is FASTA. \\
\code{--split\_pairs} & Write paired-end reads to separate \code{\_R1} and
                        \code{\_R2} files instead of interleaved. \\
\code{--maxGB <n>}   & Stop after downloading \code{n} gigabases.
                        By default the entire run is downloaded. \\
\bottomrule
\end{tabular}
\end{center}

Single-end runs are automatically exported as a single file
(\code{SRRn.fastq.gz}).  The number of SRR identifiers in a single
command is not limited.

\subsection{Streaming SRA runs on the fly}

If you have a fast connection to NCBI, you can align without storing the
raw reads locally.  Use \code{--srr} instead of \code{-i}:
\begin{lstlisting}
magic2 -x IDX --srr SRR1,SRR2,SRR3 -o out/
\end{lstlisting}

Magic2 distinguishes single-end from paired-end automatically and begins
aligning as soon as the first megabases arrive.  In practice, alignment
finishes within seconds of the download completing.

\begin{notebox}
\textbf{Note:} If the files already exist in the local \code{./SRA/}
cache directory (from a previous \code{--sraDownload}), they are used
instead of re-downloading.
\end{notebox}

To evaluate quality-control metrics across several SRA runs quickly,
without storing the reads or alignment files:
\begin{lstlisting}
magic2 -x IDX --srr SRR1,SRR2,SRR3 --no_ali -o out/
\end{lstlisting}

\subsection{Streaming and caching simultaneously}

You can align on the fly and cache the reads in one command:
\begin{lstlisting}
magic2 -x IDX.T2T --srr SRR24511885 --fastq --sraCaching -o out/ --full
# Produces: SRA/SRR24511885.sample_12.fastq.gz  (cached copy)
#           Full analysis output in out/
\end{lstlisting}

If the run is already cached, it will not be downloaded again.

\subsection{Processing multiple runs in parallel}

On large machines the index is memory-mapped and shared between
concurrent Magic2 processes.  To analyse three runs in parallel:
\begin{lstlisting}
magic2 -x IDX.T2T --srr SRR1234 --full -o out_SRR1234 &
magic2 -x IDX.T2T --srr SRR1235 --full -o out_SRR1235 &
magic2 -x IDX.T2T --srr SRR1236 --full -o out_SRR1236 &
\end{lstlisting}

\begin{notebox}
\textbf{Note:} Each parallel instance reads the shared index but uses
independent output directories.  Total RAM usage is the index size plus
per-run working memory.
\end{notebox}

% ===============================================================
\section{Building a Target Index}
% ===============================================================

\subsection{Overview}

Magic2 is a seed-and-extend aligner.  Before aligning any reads you must
build an index from your reference genome (and optionally annotation
files).  The index is written to a directory and reused for all subsequent
alignment runs against the same reference.

\subsection{The target configuration file (\code{tConfig})}

Create a plain-text configuration file with two tab-separated columns:
a single-letter type code and a file path.  Multiple lines with the
same type are allowed.

\begin{lstlisting}
G	genome.chromosomes.fasta.gz
M	genome.mitochondria.fasta.gz
R	genome.rRNA.fasta.gz
I	genome.annotation.gtf.gz
\end{lstlisting}

\textbf{Type codes:}

\begin{center}
\begin{tabular}{clp{8.5cm}}
\toprule
\textbf{Code} & \textbf{Label} & \textbf{Description} \\
\midrule
\code{G} & Genome
  & Genomic sequence.  Each chromosome may be on a separate \code{G} line. \\
\code{M} & Mitochondria
  & Mitochondrial sequence.  Reported separately because many protocols
    show high mitochondrial read fractions. \\
\code{C} & Chloroplast
  & Chloroplast sequence (plant genomes). \\
\code{R} & Ribosomal RNA
  & rRNA precursor sequences.  \textbf{Strongly recommended}: rRNA can
    constitute up to 80\% of a sample and is repeated many times in the
    genome, causing multi-mapping confusion if not handled explicitly. \\
\code{E} & Controls
  & ERCC spike-in or other sequencing controls. \\
\code{A} & Adaptors
  & Optional.  Adaptor sequences in FASTA format; if omitted, adaptors
    are detected automatically. \\
\code{B} & Bacteria
  & Potential bacterial contamination sequences. \\
\code{V} & Virus
  & Potential viral contamination sequences. \\
\code{I} & Annotation
  & Gene annotation in GTF or GFF format, or a \code{.introns} file
    (see below). \\
\bottomrule
\end{tabular}
\end{center}

\begin{notebox}
\textbf{Note:} Chromosome names in \code{G} files must match those
used in the \code{I} annotation file.
\end{notebox}

\subsubsection{Annotation file formats}

Magic2 accepts GTF, GFF, or a custom \code{.introns} format.

\textbf{GTF/GFF:} Magic2 expects at least 7 columns:
\begin{lstlisting}
chrom  method  tag  a1  a2  .  strand
\end{lstlisting}
Only \code{exon} and \code{CDS} lines are used; all others are ignored.
Coordinates are 1-based; \code{a1 < a2}; strand is \code{+} or \code{-}.

\textbf{.introns format:} If no GTF/GFF is available, provide a
3-column introns file:
\begin{lstlisting}
chrom  a1  a2
\end{lstlisting}
\code{a1} and \code{a2} give the positions of the two \texttt{G} bases
of the GT-AG motif.  On the top strand \code{a1 < a2}; on the bottom
strand \code{a1 > a2}. Coordinates are 1-based.

\subsection{Creating the index}

\textbf{From a raw genome (FASTA only):}
\begin{lstlisting}
magic2 --createIndex IDX.genome -t genome.fasta.gz
\end{lstlisting}

\textbf{From a full configuration file:}
\begin{lstlisting}
magic2 --createIndex IDX.target -T target.tConfig
\end{lstlisting}

\textbf{With custom seed parameters:}
\begin{lstlisting}
magic2 --createIndex IDX.target -T target.tConfig \
       --seedLength 18 --step 6
\end{lstlisting}

\begin{center}
\begin{tabular}{lll p{6.5cm}}
\toprule
\textbf{Option} & \textbf{Argument} & \textbf{Default} & \textbf{Description} \\
\midrule
\code{-t}            & file & ---  & Single FASTA genome file (raw genome). \\
\code{-T}            & file & ---  & Target configuration file (annotated genome). \\
\code{--seedLength}  & int  & 16 (\(\leq\)200 Mb) or 18 (human)
                     & Seed length in bases.  Not recommended to change for
                       metazoan genomes; shorter seeds may help for
                       bacterial projects. \\
\code{--step}        & int  & 6
                     & Seed stepping interval.  The best seed among the next
                       \code{step} words is selected.  Larger values (8--10)
                       reduce RAM and speed up indexing and alignments
                       at some cost to sensitivity. \\
\bottomrule
\end{tabular}
\end{center}

% ===============================================================
\section{Running an Alignment}
% ===============================================================

\subsection{Basic usage}

\begin{lstlisting}
magic2 -x IDX -i reads.fasta.gz -o out/
magic2 -x IDX -i R1.fastq.gz+R2.fastq.gz -o out/
magic2 -x IDX --srr SRR1234 -o out/
\end{lstlisting}

\subsection{Output format options}

By default, alignments are written in the standard \code{.bam} format.
You can request other formats with:
\begin{lstlisting}
magic2 -x IDX -i reads.fasta.gz -o out/ --hits --tabular --bam
\end{lstlisting}

\begin{center}
\begin{tabular}{lp{10cm}}
\toprule
\textbf{Flag} & \textbf{Description} \\
\midrule
\code{--bam}     & Standard BAM format.  Required for tools such as
                   IGV, Samtools, and most genome browsers. Default.\\
\code{--hits}    & Magic \code{.hits} format.  Describes mismatches,
                   introns, and read overhangs explicitly. \\
\code{--tabular} & Magic-BLAST tabular format.  Similar to \code{.hits}\\
\code{--no\_ali} & Suppress all alignment output.  Useful when only
                   quality metrics or expression counts are needed. \\
\bottomrule
\end{tabular}
\end{center}

Multiple format flags may be combined; one output file per format is
written to the output directory.

\subsection{Requesting the full output suite}

Adding \code{--full} produces every available output in a single pass:
\begin{lstlisting}
magic2 -x IDX.T2T --srr SRR1234 -o out/ --full
\end{lstlisting}
This is equivalent to passing \code{--genes --wiggles --wiggle\_ends}
together.  See Sections~\ref{sec:genecounts}--\ref{sec:tss} for
descriptions of each output type.

% ===============================================================
\section{Output Files}
% ===============================================================

All output files are written to the directory specified with \code{-o}.
The subsections below describe each file type.

\subsection{Alignment files}

\begin{center}
\begin{tabular}{ll p{7.5cm}}
\toprule
\textbf{File pattern} & \textbf{Format} & \textbf{Notes} \\
\midrule
\code{*.hits}         & Magic hits      & Full mismatch and intron detail. \\
\code{*.tab}          & Tabular         & Tab-separated; easy to parse. \\
\code{*.bam}          & BAM             & Standard; index with \code{samtools index}. \\
\bottomrule
\end{tabular}
\end{center}

The \code{.hits} and tabular formats provide more detail than BAM,
including explicit descriptions of mismatches, intron coordinates, and
read overhangs.  BAM is the appropriate choice when interoperating with
third-party tools.

\subsection{Intron support table}

Magic2 reports all introns observed in the alignment, classified as
\emph{known} (present in the \code{I} annotation used during index
construction) or \emph{novel}.

Output is in the \textbf{tag-sample format (TSF)}, a compact
tab-separated format in which each column corresponds to one run.
A utility command merges TSF files from multiple runs and produces a
human-readable summary table:
\begin{lstlisting}
magic2 --mergeTSF run1/introns.tsf run2/introns.tsf \
       -o merged_introns.txt
\end{lstlisting}

Each intron entry records: run, chromosome, donor position, acceptor position,
strand, GT-AG motif support, and read count per run.

\subsection{Poly-A addition sites}

An addition site is called when at least 8 terminal bases of a read
(12 if they contain a non-A/T letter) match a poly-A extension in the
3$'$ end of forward reads or a poly-T extension in the 5$'$ end of
reverse reads, and these bases do not align to the genome.  This
evidence is accumulated per genomic position and reported per run in
TSF format, analogous to the intron support table.

\subsection{Gene expression counts}
\label{sec:genecounts}

Activated by \code{--full}, \code{--wiggles}, or \code{--genes}.

Magic2 builds a 1-base-resolution coverage map internally during
alignment (at no additional I/O cost) and, if annotation was provided,
accumulates read counts per gene.

\textbf{How overlapping genes are resolved:}
The annotation is projected onto the genome and sliced into
non-overlapping regions, each classified as intergenic, intronic,
exonic (sense), exonic (antisense), or CDS.  For reads falling in
ambiguous regions (for example, two overlapping genes on opposite
strands sharing a 3$'$ UTR), attribution is deferred and distributed
between the two genes in proportion to their non-ambiguous exon support.

\textbf{Strand detection:}
The strandedness of the library is inferred from the fraction of intron
reads exhibiting GT-AG (forward) vs.\ CT-AC (reverse) motifs.  Modern
stranded protocols produce values below 1\% or above 99\%; unstranded
protocols fall in the 40--60\% range.

\textbf{Expression index:}
A preliminary expression index is computed by subtracting the
contribution of super-expressed genes (those capturing $>$1\% of all
counts), which would otherwise distort comparisons.

Counts are exported in TSF format and can be merged across runs with
the same utility used for intron tables. The merged file contains
all the counts necessary to study differential expression.

\subsection{Coverage plots}
\label{sec:wiggles}

Activated by \code{--wiggles} or \code{--full}.

Coverage plots are exported in the \textbf{BF format} (a UCSC Genome
Browser compatible wiggle format) at 10-base resolution by default.

\begin{lstlisting}
magic2 -x IDX -i reads.fasta.gz -o out/ --wiggles
magic2 -x IDX -i reads.fasta.gz -o out/ --wiggles --wiggle_step 1
\end{lstlisting}

\begin{center}
\begin{tabular}{lll p{5.5cm}}
\toprule
\textbf{Option} & \textbf{Argument} & \textbf{Default} & \textbf{Description} \\
\midrule
\code{--wiggles}      & ---  & off & Enable coverage plot output. \\
\code{--wiggle\_step} & int  & 10  & Resolution in bases.  Use 1 for
                                     single-base resolution. \\
\bottomrule
\end{tabular}
\end{center}

BF files can be loaded directly into the UCSC Genome Browser, IGV, or
converted to BigWig for large-scale display.

\subsection{Transcription start and end sites}
\label{sec:tss}

Activated by \code{--full} or \code{--wiggle\_ends}.

Magic2 exports special coverage tracks representing the first 20 bases
and final 20 bases of fully aligned reads.  Because sequencing of an
mRNA template is always oriented inward, an accumulation of read
endpoints at a particular genomic position indicates a transcription
start or termination site.  These tracks are useful for refining genome
annotations, particularly to extend or correct 5$'$ and 3$'$ UTR
boundaries.

\subsection{Quality-control metrics}

Magic2 measures and reports the following metrics for each run:

\begin{center}
\begin{longtable}{p{4.5cm} p{9cm}}
\toprule
\textbf{Metric} & \textbf{Description} \\
\midrule
\endhead
Total reads              & Total number of reads (or read pairs) in the input. \\
Aligned reads (\%)       & Fraction of reads with at least one alignment. \\
Uniquely aligned (\%)    & Fraction of reads with exactly one alignment. \\
Multi-mapping (\%)       & Fraction of reads aligning to more than one locus. \\
Unmapped (\%)            & Fraction of reads with no alignment. \\
rRNA fraction (\%)       & Fraction of reads mapping to ribosomal RNA sequences.
                           Values above 30\% may indicate inadequate rRNA depletion. \\
Mitochondrial fraction (\%) & Fraction of reads mapping to the mitochondrial sequence. \\
Adaptor rate (\%)        & Fraction of reads in which an adaptor sequence was
                           detected and clipped. \\
Library strandedness (\%)& Percentage of intron reads exhibiting GT-AG motifs
                           (vs.\ CT-AC).  Values near 0\% or 100\% indicate a
                           stranded library; values near 50\% indicate unstranded. \\
Read length distribution & Histogram of aligned read lengths (especially
                           informative for long-read data). \\
Mismatch rate (\%)       & Per-base substitution rate in aligned reads. \\
Insertion rate (\%)      & Per-base insertion rate in aligned reads. \\
Deletion rate (\%)       & Per-base deletion rate in aligned reads. \\
Soft-clip rate (\%)      & Fraction of read bases that are soft-clipped at either end. \\
Exonic fraction (\%)     & Fraction of aligned bases falling in annotated exons. \\
Intronic fraction (\%)   & Fraction of aligned bases falling in annotated introns. \\
Intergenic fraction (\%) & Fraction of aligned bases in intergenic regions. \\
\bottomrule
\end{longtable}
\end{center}

All metrics are written to a summary text file in the output directory
and also included in the TSF-format tables for multi-run comparison.

% ===============================================================
\section{Command-Line Reference}
\label{sec:cmdref}
% ===============================================================

This section lists all recognised options.  Options are grouped by
function.

\subsection{Index construction}

\begin{center}
\begin{longtable}{p{4cm} p{2.2cm} p{2cm} p{6cm}}
\toprule
\textbf{Option} & \textbf{Argument} & \textbf{Default} & \textbf{Description} \\
\midrule
\endhead
\code{--createIndex} & \textit{dir}  & required & Directory to write the index into. \\
\code{-t}            & \textit{file} & ---       & Single FASTA genome file (no annotation). \\
\code{-T}            & \textit{file} & ---       & Target configuration file (with annotation). \\
\code{--seedLength}  & \textit{int}  & 16 or 18  & Seed length in bases. \\
\code{--step}        & \textit{int}  & 6         & Seed stepping interval. \\
\bottomrule
\end{longtable}
\end{center}

\subsection{Input}

\begin{center}
\begin{longtable}{p{4cm} p{2.2cm} p{2cm} p{6cm}}
\toprule
\textbf{Option} & \textbf{Argument} & \textbf{Default} & \textbf{Description} \\
\midrule
\endhead
\code{-x}            & \textit{dir}    & required & Index directory (from \code{--createIndex}). \\
\code{-i}            & \textit{file}   & ---      & Input FASTA/FASTQ file.  Use \code{f1+f2}
                                                    for paired-end reads in two files. \\
\code{--srr}         & \textit{list}   & ---      & Comma-separated SRA run identifiers.
                                                    Streamed from NCBI or read from cache. \\
\code{--sraDownload} & \textit{list}   & ---      & Download SRA runs to \code{./SRA/} and exit. \\
\code{--fastq}       & ---             & off      & (With \code{--sraDownload}) download in
                                                    FASTQ format. \\
\code{--split\_pairs}& ---             & off      & (With \code{--sraDownload}) write pairs to
                                                    separate \code{\_R1}/\code{\_R2} files. \\
\code{--maxGB}       & \textit{int}    & unlimited& Limit download or processing to \textit{n} gigabases. \\
\code{--sraCaching}  & ---             & off      & Cache SRA data locally while aligning on the fly. \\
\bottomrule
\end{longtable}
\end{center}

\subsection{Output}

\begin{center}
\begin{longtable}{p{4cm} p{2.2cm} p{2cm} p{6cm}}
\toprule
\textbf{Option} & \textbf{Argument} & \textbf{Default} & \textbf{Description} \\
\midrule
\endhead
\code{-o}             & \textit{dir}  & required & Output directory.  Created if it does not exist. \\
\code{--hits}         & ---           & on       & Export in Magic \code{.hits} format. \\
\code{--tabular}      & ---           & off      & Export in tabular (Magic-BLAST) format. \\
\code{--bam}          & ---           & off      & Export in BAM format. \\
\code{--no\_ali}      & ---           & off      & Suppress alignment output (metrics and counts only). \\
\code{--full}         & ---           & off      & Enable all outputs: genes, wiggles, wiggle ends. \\
\code{--genes}        & ---           & off      & Export gene expression count tables. \\
\code{--wiggles}      & ---           & off      & Export coverage plots (BF/wiggle format). \\
\code{--wiggle\_ends} & ---           & off      & Export transcription-start/end coverage tracks. \\
\code{--wiggle\_step} & \textit{int}  & 10       & Coverage plot resolution in bases. \\
\bottomrule
\end{longtable}
\end{center}

% ===============================================================
\section{Worked Examples}
% ===============================================================

\subsection{Quality survey of SRA runs without storing data}

You want to assess library quality across six SRA runs before deciding
which to analyse further, without consuming disk space.

\begin{lstlisting}
magic2 -x IDX.T2T \
       --srr SRR24511880,SRR24511881,SRR24511882,\
             SRR24511883,SRR24511884,SRR24511885 \
       --no_ali -o qc_survey/
\end{lstlisting}

This streams each run from NCBI, produces quality-control metrics
(rRNA fraction, strandedness, adaptor rate, alignment rate), and
writes nothing to disk except the metric summary files.  No alignment
or coverage files are produced.

\subsection{Full analysis of a paired-end Illumina run}

\begin{lstlisting}
# Build index once (reuse for all subsequent runs)
magic2 --createIndex IDX.hg38 -T hg38.tConfig

# Analyse a local paired-end dataset
magic2 -x IDX.hg38 \
       -i sample_R1.fastq.gz+sample_R2.fastq.gz \
       -o sample_out/ \
       --bam --full

# Output directory will contain:
#   sample.bam           -- alignments (BAM)
#   introns.tsf          -- known and novel introns with support counts
#   polyA.tsf            -- poly-A addition sites
#   genes.tsf            -- gene expression counts
#   *.BF                 -- coverage wiggle tracks
#   tss.tsf / tes.tsf    -- transcription start/end site tracks
#   qc_metrics.tsf       -- quality-control summary
\end{lstlisting}

\subsection{Parallel processing of multiple SRA runs}

On a 64-core machine you can process three runs simultaneously.  The
index is memory-mapped and shared; only per-run working memory is
duplicated.

\begin{lstlisting}
magic2 -x IDX.T2T --srr SRR100001 --full -o out_SRR100001 &
magic2 -x IDX.T2T --srr SRR100002 --full -o out_SRR100002 &
magic2 -x IDX.T2T --srr SRR100003 --full -o out_SRR100003 &
wait
echo "All runs complete."
\end{lstlisting}

\subsection{Merging multi-run output tables}

After analysing several runs, merge the per-run TSF tables into a
single comparison table:
\begin{lstlisting}
magic2 --mergeTSF out_*/introns.tsf -o all_introns.txt
magic2 --mergeTSF out_*/genes.tsf   -o all_gene_counts.txt
\end{lstlisting}

% ===============================================================
\section{Troubleshooting}
% ===============================================================

\begin{center}
\begin{longtable}{p{5.5cm} p{9cm}}
\toprule
\textbf{Problem} & \textbf{Solution} \\
\midrule
\endhead

\code{magic2: command not found}
  & Ensure the directory containing the \code{magic2} binary is on your
    \code{PATH}.  Run \code{which magic2} to check. \\[4pt]

Index creation fails with out-of-memory error
  & Increase available RAM, or use a larger \code{--step} value (e.g.\
    8 or 10) to reduce index size at a some cost to sensitivity. \\[4pt]

Very low alignment rate ($<$20\%)
  & Check that the correct genome was indexed.  For RNA-seq data,
    ensure an annotation file (\code{I} line) was included in the
    tConfig, so that known splice junctions are used during alignment. \\[4pt]

Zero rRNA fraction ($>$50\%)
  & In RNA-seq experiments, this is unusual.  Ensure that an
    \code{R} line pointing to rRNA sequences was included in the
    tConfig so they are counted correctly. \\[4pt]

High rRNA fraction ($>$50\%)
  & This is a library quality issue, not an error. \\[4pt]

SRA download stalls or times out
  & Check your network connection to NCBI.  Use \code{--maxGB} to
    limit the transfer size, or pre-download with \code{--sraDownload}
    before aligning. \\[4pt]

Chromosome names don't match between genome and GTF
  & This is a fatal error reported during index creation.
    Ensure that the chromosome identifiers in your FASTA \code{G} file
    exactly match those in the \code{I} annotation file (including
    \code{chr} prefix vs.\ no prefix). \\[4pt]

BAM file not produced despite \code{--bam} flag
  & Confirm that the output directory is writable and has sufficient
    disk space.  BAM files for a human-scale genome typically require
    3--10 GB per run. \\[4pt]

Low gene count output despite \code{--full}
  & Confirm that the index was built with \code{-T} (tConfig including
    an \code{I} annotation line) rather than \code{-t} (raw genome only).
    The actual tConfig file used during index creation is copied in the
    IDX directory. Please check that file.\\

\bottomrule
\end{longtable}
\end{center}

% ===============================================================
\section*{Acknowledgements}
% ===============================================================

We thank David Landsman and Tom Madden for their support.
This work was supported by the Intramural Research Program of the
National Library of Medicine, National Institutes of Health.

% ===============================================================
\begin{thebibliography}{9}

\bibitem{cDNA}
D.~Thierry-Mieg and J.~Thierry-Mieg.
AceView: a comprehensive cDNA-supported gene and transcripts annotation.
\emph{Genome Biology}, 7(Suppl 1):S12, 2006.

\bibitem{acedb94}
R.~Durbin and J.~Thierry-Mieg.
The ACEDB genome database.
In S.~Suhai, editor, \emph{Computational Methods in Genome Research},
pages 45--55. Plenum Press, New York, 1994.

\bibitem{Code}
J.~Thierry-Mieg, D.~Thierry-Mieg, and G.~Boratyn.
Magic2 source code and documentation.
\url{https://github.com/NLM-DIR/Magic2}.

\end{thebibliography}

\end{document}
